% Options for packages loaded elsewhere
\PassOptionsToPackage{unicode}{hyperref}
\PassOptionsToPackage{hyphens}{url}
%
\documentclass[
]{article}
\usepackage{amsmath,amssymb}
\usepackage{iftex}
\ifPDFTeX
  \usepackage[T1]{fontenc}
  \usepackage[utf8]{inputenc}
  \usepackage{textcomp} % provide euro and other symbols
\else % if luatex or xetex
  \usepackage{unicode-math} % this also loads fontspec
  \defaultfontfeatures{Scale=MatchLowercase}
  \defaultfontfeatures[\rmfamily]{Ligatures=TeX,Scale=1}
\fi
\usepackage{lmodern}
\ifPDFTeX\else
  % xetex/luatex font selection
\fi
% Use upquote if available, for straight quotes in verbatim environments
\IfFileExists{upquote.sty}{\usepackage{upquote}}{}
\IfFileExists{microtype.sty}{% use microtype if available
  \usepackage[]{microtype}
  \UseMicrotypeSet[protrusion]{basicmath} % disable protrusion for tt fonts
}{}
\makeatletter
\@ifundefined{KOMAClassName}{% if non-KOMA class
  \IfFileExists{parskip.sty}{%
    \usepackage{parskip}
  }{% else
    \setlength{\parindent}{0pt}
    \setlength{\parskip}{6pt plus 2pt minus 1pt}}
}{% if KOMA class
  \KOMAoptions{parskip=half}}
\makeatother
\usepackage{xcolor}
\usepackage[margin=1in]{geometry}
\usepackage{color}
\usepackage{fancyvrb}
\newcommand{\VerbBar}{|}
\newcommand{\VERB}{\Verb[commandchars=\\\{\}]}
\DefineVerbatimEnvironment{Highlighting}{Verbatim}{commandchars=\\\{\}}
% Add ',fontsize=\small' for more characters per line
\usepackage{framed}
\definecolor{shadecolor}{RGB}{248,248,248}
\newenvironment{Shaded}{\begin{snugshade}}{\end{snugshade}}
\newcommand{\AlertTok}[1]{\textcolor[rgb]{0.94,0.16,0.16}{#1}}
\newcommand{\AnnotationTok}[1]{\textcolor[rgb]{0.56,0.35,0.01}{\textbf{\textit{#1}}}}
\newcommand{\AttributeTok}[1]{\textcolor[rgb]{0.13,0.29,0.53}{#1}}
\newcommand{\BaseNTok}[1]{\textcolor[rgb]{0.00,0.00,0.81}{#1}}
\newcommand{\BuiltInTok}[1]{#1}
\newcommand{\CharTok}[1]{\textcolor[rgb]{0.31,0.60,0.02}{#1}}
\newcommand{\CommentTok}[1]{\textcolor[rgb]{0.56,0.35,0.01}{\textit{#1}}}
\newcommand{\CommentVarTok}[1]{\textcolor[rgb]{0.56,0.35,0.01}{\textbf{\textit{#1}}}}
\newcommand{\ConstantTok}[1]{\textcolor[rgb]{0.56,0.35,0.01}{#1}}
\newcommand{\ControlFlowTok}[1]{\textcolor[rgb]{0.13,0.29,0.53}{\textbf{#1}}}
\newcommand{\DataTypeTok}[1]{\textcolor[rgb]{0.13,0.29,0.53}{#1}}
\newcommand{\DecValTok}[1]{\textcolor[rgb]{0.00,0.00,0.81}{#1}}
\newcommand{\DocumentationTok}[1]{\textcolor[rgb]{0.56,0.35,0.01}{\textbf{\textit{#1}}}}
\newcommand{\ErrorTok}[1]{\textcolor[rgb]{0.64,0.00,0.00}{\textbf{#1}}}
\newcommand{\ExtensionTok}[1]{#1}
\newcommand{\FloatTok}[1]{\textcolor[rgb]{0.00,0.00,0.81}{#1}}
\newcommand{\FunctionTok}[1]{\textcolor[rgb]{0.13,0.29,0.53}{\textbf{#1}}}
\newcommand{\ImportTok}[1]{#1}
\newcommand{\InformationTok}[1]{\textcolor[rgb]{0.56,0.35,0.01}{\textbf{\textit{#1}}}}
\newcommand{\KeywordTok}[1]{\textcolor[rgb]{0.13,0.29,0.53}{\textbf{#1}}}
\newcommand{\NormalTok}[1]{#1}
\newcommand{\OperatorTok}[1]{\textcolor[rgb]{0.81,0.36,0.00}{\textbf{#1}}}
\newcommand{\OtherTok}[1]{\textcolor[rgb]{0.56,0.35,0.01}{#1}}
\newcommand{\PreprocessorTok}[1]{\textcolor[rgb]{0.56,0.35,0.01}{\textit{#1}}}
\newcommand{\RegionMarkerTok}[1]{#1}
\newcommand{\SpecialCharTok}[1]{\textcolor[rgb]{0.81,0.36,0.00}{\textbf{#1}}}
\newcommand{\SpecialStringTok}[1]{\textcolor[rgb]{0.31,0.60,0.02}{#1}}
\newcommand{\StringTok}[1]{\textcolor[rgb]{0.31,0.60,0.02}{#1}}
\newcommand{\VariableTok}[1]{\textcolor[rgb]{0.00,0.00,0.00}{#1}}
\newcommand{\VerbatimStringTok}[1]{\textcolor[rgb]{0.31,0.60,0.02}{#1}}
\newcommand{\WarningTok}[1]{\textcolor[rgb]{0.56,0.35,0.01}{\textbf{\textit{#1}}}}
\usepackage{graphicx}
\makeatletter
\def\maxwidth{\ifdim\Gin@nat@width>\linewidth\linewidth\else\Gin@nat@width\fi}
\def\maxheight{\ifdim\Gin@nat@height>\textheight\textheight\else\Gin@nat@height\fi}
\makeatother
% Scale images if necessary, so that they will not overflow the page
% margins by default, and it is still possible to overwrite the defaults
% using explicit options in \includegraphics[width, height, ...]{}
\setkeys{Gin}{width=\maxwidth,height=\maxheight,keepaspectratio}
% Set default figure placement to htbp
\makeatletter
\def\fps@figure{htbp}
\makeatother
\setlength{\emergencystretch}{3em} % prevent overfull lines
\providecommand{\tightlist}{%
  \setlength{\itemsep}{0pt}\setlength{\parskip}{0pt}}
\setcounter{secnumdepth}{-\maxdimen} % remove section numbering
\ifLuaTeX
  \usepackage{selnolig}  % disable illegal ligatures
\fi
\usepackage{bookmark}
\IfFileExists{xurl.sty}{\usepackage{xurl}}{} % add URL line breaks if available
\urlstyle{same}
\hypersetup{
  pdftitle={EJERCICIO REGRESION},
  pdfauthor={ALONDRA DELGADO},
  hidelinks,
  pdfcreator={LaTeX via pandoc}}

\title{EJERCICIO REGRESION}
\author{ALONDRA DELGADO}
\date{2025-06-25}

\begin{document}
\maketitle

Llamamos las librerías que se van a ocupar

Cargamos el archivo desde nuestro escritorio

\begin{Shaded}
\begin{Highlighting}[]
\FunctionTok{getwd}\NormalTok{()}
\end{Highlighting}
\end{Shaded}

\begin{verbatim}
## [1] "C:/Users/emanu/OneDrive - Firedrop/Github/Public/courses/statistics/time_series/L0-Review/Review_Estadística/Scripts"
\end{verbatim}

\begin{Shaded}
\begin{Highlighting}[]
\CommentTok{\# Set working directory}
\FunctionTok{setwd}\NormalTok{(}\StringTok{"C:/Users/suscripcion35/Documents/AMAT/ST/L0{-}Review/Review\_Estadística/data"}\NormalTok{)}

\CommentTok{\# Carga el archivo "US\_change\_full.csv" dentro de la variable uschange}
\NormalTok{uschange }\OtherTok{\textless{}{-}} \FunctionTok{read.csv}\NormalTok{(}\StringTok{"../data/US\_change\_full.csv"}\NormalTok{)}
\end{Highlighting}
\end{Shaded}

Imprimimos las primeras 10 filas de nuestra base

\begin{Shaded}
\begin{Highlighting}[]
\FunctionTok{head}\NormalTok{(uschange,}\DecValTok{10}\NormalTok{)}
\end{Highlighting}
\end{Shaded}

\begin{verbatim}
##          Date Consumption     Income Production    Savings Unemployment
## 1  01/01/1970   0.6185664  1.0448013 -2.4524855  5.2990141          0.9
## 2  01/04/1970   0.4519840  1.2256472 -0.5514595  7.7898938          0.5
## 3  01/07/1970   0.8728718  1.5851538 -0.3586518  7.4039841          0.5
## 4  01/10/1970  -0.2718479 -0.2395449 -2.1856909  1.1698982          0.7
## 5  01/01/1971   1.9013450  1.9759249  1.9097644  3.5356669         -0.1
## 6  01/04/1971   0.9148773  1.4459085  0.9015695  5.8747636         -0.1
## 7  01/07/1971   0.7941096  0.5211491  0.3080307 -0.4062353          0.1
## 8  01/10/1971   1.6456332  1.1591761  2.2913621 -1.4862589          0.0
## 9  01/01/1972   1.3111897  0.4568567  4.1542950 -4.2919722         -0.2
## 10 01/04/1972   1.8857779  1.0333893  1.8886727 -4.6921964         -0.1
\end{verbatim}

Mostrar un resumen de lo que se incluye en el dataframe, en el cual
vienen variables macroeconómicas, a manera general vienen 6 variables
(fecha, consumo,ingresos, producción, ahorros y desempleo), de los
cuales existen tasas de crecimiento negativas

\begin{Shaded}
\begin{Highlighting}[]
\FunctionTok{summary}\NormalTok{(uschange)}
\end{Highlighting}
\end{Shaded}

\begin{verbatim}
##      Date            Consumption          Income          Production      
##  Length:198         Min.   :-2.2778   Min.   :-4.0844   Min.   :-6.83604  
##  Class :character   1st Qu.: 0.4175   1st Qu.: 0.3144   1st Qu.:-0.01078  
##  Mode  :character   Median : 0.7767   Median : 0.7603   Median : 0.66609  
##                     Mean   : 0.7425   Mean   : 0.7282   Mean   : 0.50708  
##                     3rd Qu.: 1.0976   3rd Qu.: 1.1602   3rd Qu.: 1.29524  
##                     Max.   : 2.3196   Max.   : 4.5219   Max.   : 4.15430  
##     Savings         Unemployment     
##  Min.   :-56.472   Min.   :-0.90000  
##  1st Qu.: -4.049   1st Qu.:-0.20000  
##  Median :  1.349   Median :-0.10000  
##  Mean   :  1.392   Mean   : 0.00101  
##  3rd Qu.:  6.341   3rd Qu.: 0.10000  
##  Max.   : 41.608   Max.   : 1.40000
\end{verbatim}

Al ser 198 filas es fácil identificar que no hay NA, sin embargo, si
fueran más datos no lo identificaríamos tan fácil, por esa razón
omitiremos los NA

\begin{Shaded}
\begin{Highlighting}[]
\NormalTok{uschange }\OtherTok{\textless{}{-}}\NormalTok{ uschange }\SpecialCharTok{\%\textgreater{}\%} \FunctionTok{na.omit}\NormalTok{()}
\FunctionTok{head}\NormalTok{(uschange,}\DecValTok{10}\NormalTok{)}
\end{Highlighting}
\end{Shaded}

\begin{verbatim}
##          Date Consumption     Income Production    Savings Unemployment
## 1  01/01/1970   0.6185664  1.0448013 -2.4524855  5.2990141          0.9
## 2  01/04/1970   0.4519840  1.2256472 -0.5514595  7.7898938          0.5
## 3  01/07/1970   0.8728718  1.5851538 -0.3586518  7.4039841          0.5
## 4  01/10/1970  -0.2718479 -0.2395449 -2.1856909  1.1698982          0.7
## 5  01/01/1971   1.9013450  1.9759249  1.9097644  3.5356669         -0.1
## 6  01/04/1971   0.9148773  1.4459085  0.9015695  5.8747636         -0.1
## 7  01/07/1971   0.7941096  0.5211491  0.3080307 -0.4062353          0.1
## 8  01/10/1971   1.6456332  1.1591761  2.2913621 -1.4862589          0.0
## 9  01/01/1972   1.3111897  0.4568567  4.1542950 -4.2919722         -0.2
## 10 01/04/1972   1.8857779  1.0333893  1.8886727 -4.6921964         -0.1
\end{verbatim}

\begin{Shaded}
\begin{Highlighting}[]
\FunctionTok{summary}\NormalTok{(uschange)}
\end{Highlighting}
\end{Shaded}

\begin{verbatim}
##      Date            Consumption          Income          Production      
##  Length:198         Min.   :-2.2778   Min.   :-4.0844   Min.   :-6.83604  
##  Class :character   1st Qu.: 0.4175   1st Qu.: 0.3144   1st Qu.:-0.01078  
##  Mode  :character   Median : 0.7767   Median : 0.7603   Median : 0.66609  
##                     Mean   : 0.7425   Mean   : 0.7282   Mean   : 0.50708  
##                     3rd Qu.: 1.0976   3rd Qu.: 1.1602   3rd Qu.: 1.29524  
##                     Max.   : 2.3196   Max.   : 4.5219   Max.   : 4.15430  
##     Savings         Unemployment     
##  Min.   :-56.472   Min.   :-0.90000  
##  1st Qu.: -4.049   1st Qu.:-0.20000  
##  Median :  1.349   Median :-0.10000  
##  Mean   :  1.392   Mean   : 0.00101  
##  3rd Qu.:  6.341   3rd Qu.: 0.10000  
##  Max.   : 41.608   Max.   : 1.40000
\end{verbatim}

La varianza individual de la variable Consumo (Consumption) se calcula
de la siguiente forma

\begin{Shaded}
\begin{Highlighting}[]
\FunctionTok{mean}\NormalTok{(uschange}\SpecialCharTok{$}\NormalTok{Consumption)}
\end{Highlighting}
\end{Shaded}

\begin{verbatim}
## [1] 0.7424821
\end{verbatim}

\begin{Shaded}
\begin{Highlighting}[]
\FunctionTok{var}\NormalTok{(uschange}\SpecialCharTok{$}\NormalTok{Consumption)}
\end{Highlighting}
\end{Shaded}

\begin{verbatim}
## [1] 0.4068692
\end{verbatim}

\begin{Shaded}
\begin{Highlighting}[]
\FunctionTok{plot}\NormalTok{(uschange}\SpecialCharTok{$}\NormalTok{Consumption)}
\end{Highlighting}
\end{Shaded}

% \includegraphics{RETO-ALO-DELGADO_files/figure-latex/unnamed-chunk-5-1.pdf}

La varianza de 0.4 indica que los datos dispersos se encuentran muy
cerca de la media

La Matriz de Varianzas/Covarianzas de las variables macroeconómicas se
ve de la siguiente forma

\begin{Shaded}
\begin{Highlighting}[]
\FunctionTok{cov}\NormalTok{(uschange[,}\DecValTok{2}\SpecialCharTok{:}\DecValTok{6}\NormalTok{])}
\end{Highlighting}
\end{Shaded}

\begin{verbatim}
##              Consumption      Income Production    Savings Unemployment
## Consumption    0.4068692  0.22037676  0.5134979  -1.958849  -0.12340208
## Income         0.2203768  0.81070679  0.3687138   7.759070  -0.07422323
## Production     0.5134979  0.36871378  2.3145044  -1.065057  -0.42889925
## Savings       -1.9588486  7.75907033 -1.0650568 143.192915   0.46572496
## Unemployment  -0.1234021 -0.07422323 -0.4288993   0.465725   0.13492283
\end{verbatim}

En la cual para la variable consumo: Cuando el consumo sube los ahorros
bajan, la producción y los ingresos suben

La matriz de correlación se ve de la siguiente forma:

\begin{Shaded}
\begin{Highlighting}[]
\FunctionTok{cor}\NormalTok{(uschange[,}\DecValTok{2}\SpecialCharTok{:}\DecValTok{6}\NormalTok{])}
\end{Highlighting}
\end{Shaded}

\begin{verbatim}
##              Consumption     Income  Production     Savings Unemployment
## Consumption    1.0000000  0.3837130  0.52915422 -0.25663310   -0.5266867
## Income         0.3837130  1.0000000  0.26917112  0.72014014   -0.2244219
## Production     0.5291542  0.2691711  1.00000000 -0.05850365   -0.7675091
## Savings       -0.2566331  0.7201401 -0.05850365  1.00000000    0.1059561
## Unemployment  -0.5266867 -0.2244219 -0.76750914  0.10595611    1.0000000
\end{verbatim}

Se observa que existe una relación alta entre los ingresos y los
ahorros, el desempleo con la producción negativamente. Por otro lado, no
existe relación entre el desempleo y ahorro, entre el ahorro y
producción

Graficandolo en un mapa de calor

\begin{Shaded}
\begin{Highlighting}[]
\FunctionTok{cor}\NormalTok{(uschange[,}\DecValTok{2}\SpecialCharTok{:}\DecValTok{6}\NormalTok{]) }\SpecialCharTok{\%\textgreater{}\%} \FunctionTok{corrplot}\NormalTok{(}\AttributeTok{method =} \StringTok{"shade"}\NormalTok{)}
\end{Highlighting}
\end{Shaded}

% \includegraphics{RETO-ALO-DELGADO_files/figure-latex/unnamed-chunk-8-1.pdf}

Ajustando un modelo de regresión lineal (simple) Y\textasciitilde X Y :
Variable respuesta Consumption X : Variable explicativa Income

Corremos lm para encontrar los coeficientes \(\beta_0\) y \(\beta_1\) en y=
\(\beta_0\)+\(\beta_1\)x Lo que está en consumo en función del ahorro

\begin{Shaded}
\begin{Highlighting}[]
\FunctionTok{plot}\NormalTok{(uschange}\SpecialCharTok{$}\NormalTok{Consumption}\SpecialCharTok{\textasciitilde{}}\NormalTok{uschange}\SpecialCharTok{$}\NormalTok{Income)}
\end{Highlighting}
\end{Shaded}

% \includegraphics{RETO-ALO-DELGADO_files/figure-latex/unnamed-chunk-9-1.pdf}

\begin{Shaded}
\begin{Highlighting}[]
\NormalTok{reg }\OtherTok{\textless{}{-}} \FunctionTok{lm}\NormalTok{(uschange}\SpecialCharTok{$}\NormalTok{Consumption}\SpecialCharTok{\textasciitilde{}}\NormalTok{uschange}\SpecialCharTok{$}\NormalTok{Income)}
\NormalTok{reg}
\end{Highlighting}
\end{Shaded}

\begin{verbatim}
## 
## Call:
## lm(formula = uschange$Consumption ~ uschange$Income)
## 
## Coefficients:
##     (Intercept)  uschange$Income  
##          0.5445           0.2718
\end{verbatim}

\begin{Shaded}
\begin{Highlighting}[]
\FunctionTok{summary}\NormalTok{(reg)}
\end{Highlighting}
\end{Shaded}

\begin{verbatim}
## 
## Call:
## lm(formula = uschange$Consumption ~ uschange$Income)
## 
## Residuals:
##      Min       1Q   Median       3Q      Max 
## -2.58236 -0.27777  0.01862  0.32330  1.42229 
## 
## Coefficients:
##                 Estimate Std. Error t value Pr(>|t|)    
## (Intercept)      0.54454    0.05403  10.079  < 2e-16 ***
## uschange$Income  0.27183    0.04673   5.817  2.4e-08 ***
## ---
## Signif. codes:  0 '***' 0.001 '**' 0.01 '*' 0.05 '.' 0.1 ' ' 1
## 
## Residual standard error: 0.5905 on 196 degrees of freedom
## Multiple R-squared:  0.1472, Adjusted R-squared:  0.1429 
## F-statistic: 33.84 on 1 and 196 DF,  p-value: 2.402e-08
\end{verbatim}

Por cada ingreso, el consumo aumenta 0.27 unidades.

p-value con los tres * significa que el intercepto tanto como nuestra
\(\beta_1\) a son significativas en nuestro modelo

Gráfica de dispersión de Consumo (eje y) VS Variable explicativa (eje x)
Incorpora una línea que indique el modelo nulo (promedio de consumo) Y
la línea de regresión

\begin{Shaded}
\begin{Highlighting}[]
\FunctionTok{plot}\NormalTok{(}\AttributeTok{x =}\NormalTok{ uschange}\SpecialCharTok{$}\NormalTok{Income,}
     \AttributeTok{y =}\NormalTok{ uschange}\SpecialCharTok{$}\NormalTok{Consumption,                    }\CommentTok{\# Coordenadas}
     \AttributeTok{col =} \FunctionTok{c}\NormalTok{(}\StringTok{"orangered1"}\NormalTok{),               }\CommentTok{\# De qué color (puede ser más de uno e incluso ponerle "colors()")}
     \AttributeTok{pch =} \DecValTok{18}\NormalTok{,                            }\CommentTok{\# Tipo de punto que se va a utilizar}
     \AttributeTok{main =} \StringTok{"Ingresos vs Consumo"}\NormalTok{,     }\CommentTok{\# Título del gráfico}
     \AttributeTok{xlab =} \StringTok{"Ingreso per cápita (L)"}\NormalTok{, }\CommentTok{\# Nombre del eje x}
     \AttributeTok{ylab =} \StringTok{"Consumo"}\NormalTok{)                    }\CommentTok{\# Nombre del eje y}

\CommentTok{\# abline pinta sobre el gráfico actual una recta del tipo y = a + bx}

\CommentTok{\# Pintemos el MODELO NULO, es decir la media de y}
\FunctionTok{abline}\NormalTok{(}\AttributeTok{a =} \FunctionTok{mean}\NormalTok{(uschange}\SpecialCharTok{$}\NormalTok{Consumption), }\AttributeTok{b =} \DecValTok{0}\NormalTok{, }\AttributeTok{col =} \StringTok{"blue"}\NormalTok{, }\AttributeTok{lwd =} \DecValTok{2}\NormalTok{)}

\CommentTok{\# Intentos dse encontrar una recta que se ajuste a los datos}
\FunctionTok{abline}\NormalTok{(}\AttributeTok{a =} \FloatTok{0.3}\NormalTok{, }\AttributeTok{b =} \FloatTok{0.5}\NormalTok{, }\AttributeTok{col =} \StringTok{"gray50"}\NormalTok{, }\AttributeTok{lwd =} \DecValTok{1}\NormalTok{)}
\end{Highlighting}
\end{Shaded}

% \includegraphics{RETO-ALO-DELGADO_files/figure-latex/unnamed-chunk-10-1.pdf}

Guardar las predicciones y residuales dentro del data frame Y elaborar
una regresión de residuales \textasciitilde{} predicciones

\begin{Shaded}
\begin{Highlighting}[]
\NormalTok{uschange}\SpecialCharTok{$}\NormalTok{pred }\OtherTok{\textless{}{-}}\NormalTok{ reg}\SpecialCharTok{$}\NormalTok{fitted.values}
\NormalTok{uschange}\SpecialCharTok{$}\NormalTok{res }\OtherTok{\textless{}{-}}\NormalTok{ reg}\SpecialCharTok{$}\NormalTok{residuals}

\CommentTok{\# Regresión de residuales}
\NormalTok{reg\_res }\OtherTok{\textless{}{-}} \FunctionTok{lm}\NormalTok{(uschange}\SpecialCharTok{$}\NormalTok{res}\SpecialCharTok{\textasciitilde{}}\NormalTok{uschange}\SpecialCharTok{$}\NormalTok{pred) }\CommentTok{\# "\textasciitilde{}" Regresión de los residuales con base en predicciones}
\NormalTok{reg\_res}
\end{Highlighting}
\end{Shaded}

\begin{verbatim}
## 
## Call:
## lm(formula = uschange$res ~ uschange$pred)
## 
## Coefficients:
##   (Intercept)  uschange$pred  
##     1.918e-16     -2.424e-16
\end{verbatim}

\begin{Shaded}
\begin{Highlighting}[]
\FunctionTok{summary}\NormalTok{(reg\_res)}
\end{Highlighting}
\end{Shaded}

\begin{verbatim}
## 
## Call:
## lm(formula = uschange$res ~ uschange$pred)
## 
## Residuals:
##      Min       1Q   Median       3Q      Max 
## -2.58236 -0.27777  0.01862  0.32330  1.42229 
## 
## Coefficients:
##                 Estimate Std. Error t value Pr(>|t|)
## (Intercept)    1.918e-16  1.344e-01       0        1
## uschange$pred -2.424e-16  1.719e-01       0        1
## 
## Residual standard error: 0.5905 on 196 degrees of freedom
## Multiple R-squared:  2.24e-32,   Adjusted R-squared:  -0.005102 
## F-statistic: 4.39e-30 on 1 and 196 DF,  p-value: 1
\end{verbatim}

Elaborar el gráfico de residuales VS predicciones

\begin{Shaded}
\begin{Highlighting}[]
\NormalTok{uschange }\SpecialCharTok{\%\textgreater{}\%} \FunctionTok{ggplot}\NormalTok{(}\FunctionTok{aes}\NormalTok{(}\AttributeTok{x =}\NormalTok{ uschange}\SpecialCharTok{$}\NormalTok{pred, }\AttributeTok{y =}\NormalTok{ uschange}\SpecialCharTok{$}\NormalTok{res)) }\SpecialCharTok{+} \FunctionTok{geom\_point}\NormalTok{(}\AttributeTok{alpha =} \FloatTok{0.6}\NormalTok{, }\AttributeTok{color =} \StringTok{"\#001F82"}\NormalTok{) }\SpecialCharTok{+} 
  \FunctionTok{labs}\NormalTok{(}\AttributeTok{title =} \StringTok{"Predicciones VS Residuales"}\NormalTok{,}
       \AttributeTok{subtitle =} \StringTok{"Gráfico de dispersión"}\NormalTok{,}
       \AttributeTok{x =} \StringTok{"Predicción Ingreso"}\NormalTok{,}
       \AttributeTok{y =} \StringTok{"Residuales"}\NormalTok{) }\SpecialCharTok{+}
  \FunctionTok{theme\_minimal}\NormalTok{() }\SpecialCharTok{+}
  \FunctionTok{theme}\NormalTok{(}
    \AttributeTok{plot.title =} \FunctionTok{element\_text}\NormalTok{(}\AttributeTok{color =} \StringTok{"\#0099F8"}\NormalTok{,}
                              \AttributeTok{size =} \DecValTok{17}\NormalTok{,}
                              \AttributeTok{face =} \StringTok{"bold"}\NormalTok{),}
    \AttributeTok{plot.subtitle =} \FunctionTok{element\_text}\NormalTok{(}\AttributeTok{color =} \StringTok{"\#969696"}\NormalTok{, }\AttributeTok{size =} \DecValTok{13}\NormalTok{, }\AttributeTok{face =} \StringTok{"italic"}\NormalTok{),}
    \AttributeTok{axis.title =} \FunctionTok{element\_text}\NormalTok{(}\AttributeTok{color =} \StringTok{"\#969696"}\NormalTok{,}
                              \AttributeTok{size =} \DecValTok{10}\NormalTok{,}
                              \AttributeTok{face =} \StringTok{"bold"}\NormalTok{),}
    \AttributeTok{axis.text =} \FunctionTok{element\_text}\NormalTok{(}\AttributeTok{color =} \StringTok{"\#969696"}\NormalTok{, }\AttributeTok{size =} \DecValTok{10}\NormalTok{),}
    \AttributeTok{axis.line =} \FunctionTok{element\_line}\NormalTok{(}\AttributeTok{color =} \StringTok{"\#969696"}\NormalTok{)}
\NormalTok{  )}
\end{Highlighting}
\end{Shaded}

\begin{verbatim}
## Warning: Use of `uschange$pred` is discouraged.
## i Use `pred` instead.
\end{verbatim}

\begin{verbatim}
## Warning: Use of `uschange$res` is discouraged.
## i Use `res` instead.
\end{verbatim}

% \includegraphics{RETO-ALO-DELGADO_files/figure-latex/unnamed-chunk-12-1.pdf}
Hagamos nuevamente el gráfico con la línea de regresión,

\begin{Shaded}
\begin{Highlighting}[]
\NormalTok{uschange }\SpecialCharTok{\%\textgreater{}\%} \FunctionTok{ggplot}\NormalTok{(}\FunctionTok{aes}\NormalTok{(}\AttributeTok{x =}\NormalTok{ pred, }\AttributeTok{y =}\NormalTok{ res)) }\SpecialCharTok{+} \FunctionTok{geom\_point}\NormalTok{(}\AttributeTok{alpha =} \FloatTok{0.6}\NormalTok{, }\AttributeTok{color =} \StringTok{"\#001F82"}\NormalTok{) }\SpecialCharTok{+} 
  \FunctionTok{geom\_abline}\NormalTok{(}\AttributeTok{intercept =} \FunctionTok{coef}\NormalTok{(reg\_res)[}\DecValTok{1}\NormalTok{], }\AttributeTok{slope =} \FunctionTok{coef}\NormalTok{(reg\_res)[}\DecValTok{2}\NormalTok{], }\AttributeTok{color =} \StringTok{"\#0099F8"}\NormalTok{, }\AttributeTok{size =} \FloatTok{1.5}\NormalTok{)}\SpecialCharTok{+}
  \FunctionTok{geom\_text}\NormalTok{(}
    \AttributeTok{label=}\NormalTok{ uschange}\SpecialCharTok{$}\NormalTok{Date,}
    \AttributeTok{nudge\_x =} \DecValTok{0}\NormalTok{, }\AttributeTok{nudge\_y =} \DecValTok{2}\NormalTok{,}
    \AttributeTok{check\_overlap =}\NormalTok{ T}
\NormalTok{  ) }\SpecialCharTok{+}
  \FunctionTok{geom\_label}\NormalTok{(}
    \AttributeTok{data =}\NormalTok{ uschange }\SpecialCharTok{\%\textgreater{}\%} \FunctionTok{filter}\NormalTok{(pred }\SpecialCharTok{\textless{}} \DecValTok{1}\NormalTok{), }\CommentTok{\# Filtramos datos}
    \FunctionTok{aes}\NormalTok{(}\AttributeTok{label =}\NormalTok{ Date,}
        \AttributeTok{x =}\NormalTok{ pred,}
        \AttributeTok{y =}\NormalTok{ res),}
    \AttributeTok{nudge\_x =} \DecValTok{0}\NormalTok{, }\AttributeTok{nudge\_y =} \DecValTok{0}\NormalTok{,}
    \AttributeTok{label.size =} \FloatTok{0.3}\NormalTok{,}
    \AttributeTok{label.padding =} \FunctionTok{unit}\NormalTok{(}\FloatTok{0.15}\NormalTok{, }\StringTok{"lines"}\NormalTok{),}
    \AttributeTok{fill=}\StringTok{"lightblue"}\NormalTok{) }\SpecialCharTok{+}
  \FunctionTok{labs}\NormalTok{(}\AttributeTok{title =} \StringTok{"Predicciones VS Residuales"}\NormalTok{,}
     \AttributeTok{subtitle =} \StringTok{"Gráfico de dispersión"}\NormalTok{,}
     \AttributeTok{x =} \StringTok{"Predicción consumo"}\NormalTok{,}
     \AttributeTok{y =} \StringTok{"Residuales"}\NormalTok{) }\SpecialCharTok{+}
  \FunctionTok{theme\_minimal}\NormalTok{() }\SpecialCharTok{+}
  \FunctionTok{theme}\NormalTok{(}
    \AttributeTok{plot.title =} \FunctionTok{element\_text}\NormalTok{(}\AttributeTok{color =} \StringTok{"\#0099F8"}\NormalTok{,}
                              \AttributeTok{size =} \DecValTok{17}\NormalTok{,}
                              \AttributeTok{face =} \StringTok{"bold"}\NormalTok{),}
    \AttributeTok{plot.subtitle =} \FunctionTok{element\_text}\NormalTok{(}\AttributeTok{color =} \StringTok{"\#969696"}\NormalTok{, }\AttributeTok{size =} \DecValTok{13}\NormalTok{, }\AttributeTok{face =} \StringTok{"italic"}\NormalTok{),}
    \AttributeTok{axis.title =} \FunctionTok{element\_text}\NormalTok{(}\AttributeTok{color =} \StringTok{"\#969696"}\NormalTok{,}
                              \AttributeTok{size =} \DecValTok{10}\NormalTok{,}
                              \AttributeTok{face =} \StringTok{"bold"}\NormalTok{),}
    \AttributeTok{axis.text =} \FunctionTok{element\_text}\NormalTok{(}\AttributeTok{color =} \StringTok{"\#969696"}\NormalTok{, }\AttributeTok{size =} \DecValTok{10}\NormalTok{),}
    \AttributeTok{axis.line =} \FunctionTok{element\_line}\NormalTok{(}\AttributeTok{color =} \StringTok{"\#969696"}\NormalTok{)}
\NormalTok{  )}
\end{Highlighting}
\end{Shaded}

\begin{verbatim}
## Warning: Using `size` aesthetic for lines was deprecated in ggplot2 3.4.0.
## i Please use `linewidth` instead.
## This warning is displayed once every 8 hours.
## Call `lifecycle::last_lifecycle_warnings()` to see where this warning was
## generated.
\end{verbatim}

% \includegraphics{RETO-ALO-DELGADO_files/figure-latex/unnamed-chunk-13-1.pdf}

No se alcanzan a percibir muy bien los datos con la fecha dada la
cantidad de datos.

Regresión múltiple Elabora un modelo de regresión múltiple Variable (Y)
dependiente sea Consumo Variables explicativas
(Ingreso,ahorro,produccion, desempleo)

\begin{Shaded}
\begin{Highlighting}[]
\CommentTok{\#Modelo 1 (ingreso, ahorro)}
\NormalTok{reg\_mult }\OtherTok{\textless{}{-}} \FunctionTok{lm}\NormalTok{(uschange}\SpecialCharTok{$}\NormalTok{Consumption}\SpecialCharTok{\textasciitilde{}}\NormalTok{uschange}\SpecialCharTok{$}\NormalTok{Income }\SpecialCharTok{+}\NormalTok{ uschange}\SpecialCharTok{$}\NormalTok{Savings) }
\NormalTok{reg\_mult}
\end{Highlighting}
\end{Shaded}

\begin{verbatim}
## 
## Call:
## lm(formula = uschange$Consumption ~ uschange$Income + uschange$Savings)
## 
## Coefficients:
##      (Intercept)   uschange$Income  uschange$Savings  
##          0.21543           0.83664          -0.05901
\end{verbatim}

\begin{Shaded}
\begin{Highlighting}[]
\FunctionTok{summary}\NormalTok{(reg\_mult)}
\end{Highlighting}
\end{Shaded}

\begin{verbatim}
## 
## Call:
## lm(formula = uschange$Consumption ~ uschange$Income + uschange$Savings)
## 
## Residuals:
##      Min       1Q   Median       3Q      Max 
## -0.96934 -0.14829 -0.01167  0.14240  1.41274 
## 
## Coefficients:
##                  Estimate Std. Error t value Pr(>|t|)    
## (Intercept)       0.21543    0.03393   6.349 1.49e-09 ***
## uschange$Income   0.83664    0.03748  22.324  < 2e-16 ***
## uschange$Savings -0.05901    0.00282 -20.927  < 2e-16 ***
## ---
## Signif. codes:  0 '***' 0.001 '**' 0.01 '*' 0.05 '.' 0.1 ' ' 1
## 
## Residual standard error: 0.3286 on 195 degrees of freedom
## Multiple R-squared:  0.7373, Adjusted R-squared:  0.7346 
## F-statistic: 273.6 on 2 and 195 DF,  p-value: < 2.2e-16
\end{verbatim}

\begin{Shaded}
\begin{Highlighting}[]
\CommentTok{\#Modelo 2 (Ingreso,ahorro,produccion,desempleo)}
\CommentTok{\#reg\_mult2 \textless{}{-} lm(uschange$Consumption\textasciitilde{}uschange$Income + uschange$Savings + uschange$Production + uschange$Unemployment) }
\CommentTok{\#reg\_mult2}
\CommentTok{\#summary(reg\_mult2)}
\end{Highlighting}
\end{Shaded}

Para el modelo 1: Por cada unidad adicional de ingreso, el consumo
aumenta en aproximadamente 0.84 unidades. Por cada unidad adicional de
ahorro, el consumo disminuye en aproximadamente 0.059 unidades

Ambas tienen sentido dado que entre más ganemos, mas consumimos y por
otro lado, entre más ahorremos, menos tenemos para consumir, por lo que
disminuye el consumo.

En cuanto a los niveles de significancia ambas variables son
significativas para el modelo.

Guardando las predicciones y residuales Y elabora una regresión de
residuales \textasciitilde{} predicciones Para el modelo múltiple

\begin{Shaded}
\begin{Highlighting}[]
\NormalTok{uschange}\SpecialCharTok{$}\NormalTok{pred\_mult }\OtherTok{\textless{}{-}}\NormalTok{ reg\_mult}\SpecialCharTok{$}\NormalTok{fitted.values}
\NormalTok{uschange}\SpecialCharTok{$}\NormalTok{res\_mult }\OtherTok{\textless{}{-}}\NormalTok{ reg\_mult}\SpecialCharTok{$}\NormalTok{residuals}

\CommentTok{\# Regresión de residuales}
\NormalTok{reg\_res\_mult }\OtherTok{\textless{}{-}} \FunctionTok{lm}\NormalTok{(uschange}\SpecialCharTok{$}\NormalTok{res\_mult}\SpecialCharTok{\textasciitilde{}}\NormalTok{uschange}\SpecialCharTok{$}\NormalTok{pred\_mult) }\CommentTok{\# "\textasciitilde{}" Regresión de los residuales con base en predicciones}
\NormalTok{reg\_res\_mult}
\end{Highlighting}
\end{Shaded}

\begin{verbatim}
## 
## Call:
## lm(formula = uschange$res_mult ~ uschange$pred_mult)
## 
## Coefficients:
##        (Intercept)  uschange$pred_mult  
##          3.389e-18          -7.221e-18
\end{verbatim}

\begin{Shaded}
\begin{Highlighting}[]
\FunctionTok{summary}\NormalTok{(reg\_res\_mult)}
\end{Highlighting}
\end{Shaded}

\begin{verbatim}
## 
## Call:
## lm(formula = uschange$res_mult ~ uschange$pred_mult)
## 
## Residuals:
##      Min       1Q   Median       3Q      Max 
## -0.96934 -0.14829 -0.01167  0.14240  1.41274 
## 
## Coefficients:
##                      Estimate Std. Error t value Pr(>|t|)
## (Intercept)         3.389e-18  3.930e-02       0        1
## uschange$pred_mult -7.221e-18  4.264e-02       0        1
## 
## Residual standard error: 0.3278 on 196 degrees of freedom
## Multiple R-squared:  9.853e-33,  Adjusted R-squared:  -0.005102 
## F-statistic: 1.931e-30 on 1 and 196 DF,  p-value: 1
\end{verbatim}

Gráfico de residuales contra predicciones para el modelo múltiple

\begin{Shaded}
\begin{Highlighting}[]
\NormalTok{uschange }\SpecialCharTok{\%\textgreater{}\%} \FunctionTok{ggplot}\NormalTok{(}\FunctionTok{aes}\NormalTok{(}\AttributeTok{x =}\NormalTok{ uschange}\SpecialCharTok{$}\NormalTok{pred\_mult, }\AttributeTok{y =}\NormalTok{ uschange}\SpecialCharTok{$}\NormalTok{res\_mult)) }\SpecialCharTok{+} \FunctionTok{geom\_point}\NormalTok{(}\AttributeTok{alpha =} \FloatTok{0.6}\NormalTok{, }\AttributeTok{color =} \StringTok{"\#001F82"}\NormalTok{) }\SpecialCharTok{+} 
  \FunctionTok{geom\_abline}\NormalTok{(}\AttributeTok{intercept =} \FunctionTok{coef}\NormalTok{(reg\_res\_mult)[}\DecValTok{1}\NormalTok{], }\AttributeTok{slope =} \FunctionTok{coef}\NormalTok{(reg\_res\_mult)[}\DecValTok{2}\NormalTok{], }\AttributeTok{color =} \StringTok{"\#0099F8"}\NormalTok{, }\AttributeTok{size =} \FloatTok{1.5}\NormalTok{)}\SpecialCharTok{+}
  \FunctionTok{labs}\NormalTok{(}\AttributeTok{title =} \StringTok{"Predicciones VS Residuales"}\NormalTok{,}
       \AttributeTok{subtitle =} \StringTok{"Consumo \textasciitilde{} Ingreso + Ahorro"}\NormalTok{,}
       \AttributeTok{x =} \StringTok{"Predicción consumo"}\NormalTok{,}
       \AttributeTok{y =} \StringTok{"Residuales"}\NormalTok{) }\SpecialCharTok{+}
  \FunctionTok{theme\_minimal}\NormalTok{() }\SpecialCharTok{+}
  \FunctionTok{theme}\NormalTok{(}
    \AttributeTok{plot.title =} \FunctionTok{element\_text}\NormalTok{(}\AttributeTok{color =} \StringTok{"\#0099F8"}\NormalTok{,}
                              \AttributeTok{size =} \DecValTok{17}\NormalTok{,}
                              \AttributeTok{face =} \StringTok{"bold"}\NormalTok{),}
    \AttributeTok{plot.subtitle =} \FunctionTok{element\_text}\NormalTok{(}\AttributeTok{color =} \StringTok{"\#969696"}\NormalTok{, }\AttributeTok{size =} \DecValTok{13}\NormalTok{, }\AttributeTok{face =} \StringTok{"italic"}\NormalTok{),}
    \AttributeTok{axis.title =} \FunctionTok{element\_text}\NormalTok{(}\AttributeTok{color =} \StringTok{"\#969696"}\NormalTok{,}
                              \AttributeTok{size =} \DecValTok{10}\NormalTok{,}
                              \AttributeTok{face =} \StringTok{"bold"}\NormalTok{),}
    \AttributeTok{axis.text =} \FunctionTok{element\_text}\NormalTok{(}\AttributeTok{color =} \StringTok{"\#969696"}\NormalTok{, }\AttributeTok{size =} \DecValTok{10}\NormalTok{),}
    \AttributeTok{axis.line =} \FunctionTok{element\_line}\NormalTok{(}\AttributeTok{color =} \StringTok{"\#969696"}\NormalTok{)}
\NormalTok{  )}
\end{Highlighting}
\end{Shaded}

\begin{verbatim}
## Warning: Use of `uschange$pred_mult` is discouraged.
## i Use `pred_mult` instead.
\end{verbatim}

\begin{verbatim}
## Warning: Use of `uschange$res_mult` is discouraged.
## i Use `res_mult` instead.
\end{verbatim}

% \includegraphics{RETO-ALO-DELGADO_files/figure-latex/unnamed-chunk-16-1.pdf}

Comparación de ambos modelos

\begin{Shaded}
\begin{Highlighting}[]
\FunctionTok{glance}\NormalTok{(reg\_res\_mult)}
\end{Highlighting}
\end{Shaded}

\begin{verbatim}
## # A tibble: 1 x 12
##   r.squared adj.r.squared sigma statistic p.value    df logLik   AIC   BIC
##       <dbl>         <dbl> <dbl>     <dbl>   <dbl> <dbl>  <dbl> <dbl> <dbl>
## 1  9.85e-33      -0.00510 0.328  1.93e-30    1.00     1  -59.1  124.  134.
## # i 3 more variables: deviance <dbl>, df.residual <int>, nobs <int>
\end{verbatim}

\begin{Shaded}
\begin{Highlighting}[]
\FunctionTok{glance}\NormalTok{(reg\_res)}
\end{Highlighting}
\end{Shaded}

\begin{verbatim}
## # A tibble: 1 x 12
##   r.squared adj.r.squared sigma statistic p.value    df logLik   AIC   BIC
##       <dbl>         <dbl> <dbl>     <dbl>   <dbl> <dbl>  <dbl> <dbl> <dbl>
## 1  2.24e-32      -0.00510 0.591  4.39e-30    1.00     1  -176.  357.  367.
## # i 3 more variables: deviance <dbl>, df.residual <int>, nobs <int>
\end{verbatim}

Con base en el criterio de AIC el mejor modelo sería el segundo modelo
con las variables ingreso y ahorro.

\end{document}
